\documentclass[a4paper,oneside,12pt]{book}

%----------------------------------------------------------------------------------------
%	README!
%   Welcome. It's worth having a read through this file
%   to set up the broad parameters, such as the name of
%   the degree, the school/department, the type of work
%   (dissertation/Final Year Project/report, etc. as well
%   as your own details.
%----------------------------------------------------------------------------------------

%----------------------------------------------------------------------------------------
%	COVER PAGE
%   The cover page is laid out in title/title.tex. You can choose a colour
%   or black and white logo
%----------------------------------------------------------------------------------------

%----------------------------------------------------------------------------------------
%	THESIS INFORMATION
%   Put title, author name, supervisor name, degree, type of work, school, department in here
%   It will be used for the title page and for the embedded PDF information
%----------------------------------------------------------------------------------------

\newcommand{\thesistitle}{Consumer Perception of AI-Mediated Communication in Digital Marketing Contexts: A Thematic Literature Review} % Your thesis title, this is used in the title and abstract
\newcommand{\degree}{M.Sc. in Interactive Digital Media} % Replace with your degree name, this is used in the title page and abstract
\newcommand{\typeofthesis}{dissertation} % dissertation, Final Year Project, report, etc.
\newcommand{\authorname}{Pojui Hsiao} % Your name, this is used in the title page and PDF stuff
%% Do not put your Student ID in the document, as TCD will not publish
%% documents that contain both your name and your Student ID.
\newcommand{\supervisor}{Carol O'Sullivan} % replace with the name of your supervisor
%\newcommand{\cosupervisor}{Dr Alex Lee} % replace with the name of your co-supervisor if you have one
\newcommand{\keywords}{artificial intelligence, generative artificial intelligence, AI-mediated communication, digital marketing, consumer perception, trust, authenticity} % Replace with keywords for your thesis
\newcommand{\school}{\href{https://www.scss.tcd.ie/}{School of Computer Science and Statistics}} % Your school's name and URL, this is used in the title page
%Edited by HS for engineering

%% Comment out the next line if you don't want a department to appear
%\newcommand{\department}{\href{http://researchgroup.university.com}{Department Name}} % Your research group's name and URL, this is used in the title page


%% Language and font encodings
\usepackage[T1]{fontenc} 
\usepackage[utf8]{inputenc}
\usepackage[english]{babel}

%% Bibliographical stuff
\usepackage[]{cite}
%% Document size
% include showframe as an option if you want to see the boxes
\usepackage[a4paper,top=2.54cm,bottom=2.54cm,left=2.54cm,right=2.54cm,headheight=16pt]{geometry}
\setlength{\marginparwidth}{2cm}
%% Useful packages
\usepackage{amsmath}
\usepackage[autostyle=true]{csquotes} % Required to generate language-dependent quotes in the bibliography
\usepackage[pdftex]{graphicx}
\usepackage[colorinlistoftodos]{todonotes}
\usepackage[colorlinks=true, allcolors=black]{hyperref}
\usepackage{hyperxmp}
\usepackage{caption} % if no caption, no colon
\usepackage{sfmath} %use sans-serif in the maths sections too
\usepackage[parfill]{parskip}    % Begin paragraphs with an empty line rather than an indent
\usepackage{setspace} % to permit one-and-a-half or double spacing
\usepackage{enumerate} % fancy enumerations like (i) (ii) or (a) (b) and suchlike
\usepackage{booktabs} % To thicken table lines
\usepackage{fancyhdr}
\usepackage{xcolor} % to get TCD colour on headings
\numberwithin{equation}{chapter} %HS edit for (chapter.equation)
\pagestyle{plain} % Embrace simplicity!

\definecolor{tcd_blue}{RGB}{5, 105, 185}

%% It's personal taste but...
%% Uncomment the following block if you want your name and ID at the top of
%% (almost) every page.

%\pagestyle{fancy}
%\fancyhf{} % sets both header and footer to nothing
%\renewcommand{\headrulewidth}{0pt}
%\cfoot{\thepage}
%\ifdefined\authorid
%\chead{\it \authorname\ (\authorid)}
%\else
%\chead{\it \authorname}
%\fi
%% End of block

%% It is good practise to make your font sans-serif to improve the accessibility of your document.  Comment out the following line to disable it (but you really should not)
\renewcommand{\familydefault}{\sfdefault} %use the sans-serif font as default

%% If you insist on not using sans-serif (please don't), consider using Palatino instead of the LaTeX standard
%\usepackage{mathpazo} % Use the Palatino font by default if you prefer it to Computer Modern


%% Format Chapter headings appropriately
\usepackage{titlesec}
\titleformat{\chapter}[hang]{\normalfont\huge\bfseries\color{tcd_blue}}{\thechapter}{1cm}{}{}

\title{\thesistitle}
\author{\authorname}


\hypersetup{
   pdftitle=\thesistitle, % Set the PDF's title to your title
   pdfauthor=\authorname, % Set the PDF's author to your name
   pdfkeywords=\keywords, % Set the PDF's keywords to your keywords
   pdfsubject=\degree, % Set the PDF's keywords to your keywords
   pdfinfo={
     pdfsupervisor=\supervisor, % Set the PDF's supervisor to your supervisor
     %pdfcosupervisor=\cosupervisor, % Set the PDF's cosupervisor to your cosupervisor if using
   }
}


\frontmatter
\begin{document}
\input{title/title.tex}
\pagenumbering{roman}
\section*{\Huge\textcolor{tcd_blue}{Declaration}}
\vspace{1cm}
I hereby declare that this \typeofthesis\ is entirely my own work and that it has not been submitted as an exercise for a degree at this or any other university.

\vspace{1cm}
I have read and I understand the plagiarism provisions in the General Regulations of the University Calendar for the current year, found at \url{http://www.tcd.ie/calendar}.
\vspace{1cm}

I have completed the Online Tutorial on avoiding plagiarism `Ready Steady Write', located at \url{http://tcd-ie.libguides.com/plagiarism/ready-steady-write}.
\vspace{1cm}

I consent / do not consent to the examiner retaining a copy of the thesis beyond the examining period, should they so wish (EU GDPR May 2018).
\vspace{1cm}

I agree that this thesis will not be publicly available, but will be available to TCD staff and students in the University’s open access institutional repository on the Trinity domain only, subject to Irish Copyright Legislation and Trinity College Library conditions of use and acknowledgement.  \textbf{Please consult with your supervisor on this last item before agreeing, and delete if you do not consent}
\vspace{3cm}

Signed:~\rule{5cm}{0.3pt}\hfill Date:~\rule{5cm}{0.3pt}

\chapter*{Abstract}
This dissertation examines how artificial intelligence (AI) and generative artificial intelligence (GenAI) shape consumer perception in AI-mediated communication across digital marketing contexts. As AI becomes more visible in personalised advertising, recommendation systems, service interactions, and content creation, consumer response is shaped not only by efficiency and convenience but also by concerns about trust, authenticity, privacy, and behavioural influence. Although the literature on AI in marketing has expanded rapidly, it remains fragmented across applications and often isolates individual constructs rather than explaining how they interact.

To address this gap, the study adopts a qualitative literature review using thematic analysis. Peer-reviewed academic studies on AI, GenAI, and consumer perception in digital communication and marketing settings were selected and analysed to identify recurring themes and conceptual tensions. The review focuses on how consumers interpret AI-mediated brand communication and how these interpretations affect attitudes and behavioural intentions.

The findings show that consumer perception of AI-mediated communication is neither uniformly positive nor uniformly negative. Trust emerges as a central mediating factor, while authenticity becomes especially important when AI is involved in emotionally expressive or creative communication. Personalisation can improve perceived relevance and usefulness, but its benefits are often constrained by privacy concerns and perceived intrusiveness. Behavioural outcomes such as purchase intention, engagement, and donation intention are therefore shaped by the interaction of several perceptual mechanisms rather than by AI adoption alone.

The dissertation concludes that AI-mediated communication should be understood as a consumer-facing relational issue rather than only a technological or managerial one. It contributes to the literature by offering a consumer-centred synthesis of the main themes that shape responses to AI and GenAI in digital brand communication.

\chapter*{Lay Abstract}
This dissertation looks at how people feel when brands and businesses use artificial intelligence to communicate with them in digital media. Today, AI is used to recommend products, personalise advertisements, answer customer questions, and even create promotional text or images. While these tools can make communication faster and more relevant, they can also make people uneasy. Some people may find AI useful and convenient, while others may feel that it is intrusive, less trustworthy, or less genuine than human communication.

The study does not collect survey or interview data. Instead, it reviews and analyses existing academic research on AI in digital marketing and brand communication. By comparing published studies, the dissertation identifies the main patterns in how consumers react to AI-generated or AI-supported communication.

The review shows that consumer reactions depend on more than just whether AI is effective. People are more likely to respond positively when AI-mediated communication feels helpful, relevant, and easy to use. However, negative reactions become more likely when AI reduces a sense of authenticity, raises privacy concerns, or makes consumers feel that communication is too automated. Overall, the dissertation argues that successful use of AI depends on how people interpret its presence in digital communication, not just on what the technology can do.

\newpage
\onehalfspacing\raggedright %\raggedright turns off justification and hyphenation

\section*{\Huge\textcolor{tcd_blue}{Acknowledgements}}
I would like to express my sincere gratitude to my supervisor for their guidance, encouragement, and feedback throughout this dissertation. Their advice helped me refine the focus of the study and develop a clearer understanding of the literature.

I am also grateful to the lecturers and staff of the School of Computer Science and Statistics for their support during my studies in Interactive Digital Media. Finally, I would like to thank my family and friends for their patience and encouragement throughout the research and writing process.
\tableofcontents
\listoffigures
\listoftables

\mainmatter
% maintaining separate .tex files for each chapter is good practice
\chapter{Introduction}

\section{Research Background}
Artificial intelligence (AI) is reshaping the interfaces through which brands and consumers meet in digital environments. Across social platforms, e-commerce systems, search interfaces, chatbots, and content feeds, AI increasingly helps determine how messages are produced, personalised, delivered, and interpreted. This matters not only because AI changes marketing efficiency, but because it changes the communicative form of digital brand interaction itself. Recent scholarship suggests that AI is no longer a peripheral back-end tool used only for optimisation; it increasingly participates in the design of consumer-facing communication, interaction, and media experience \cite{haleem2022ai,kumar2024aipowered,jain2024twodecades}.

Recent industry evidence supports the scale of this transition. McKinsey reports that 88\% of organisations now use AI regularly in at least one business function, with marketing and sales among the most common deployment areas, yet most firms remain in experimentation or piloting rather than full-scale transformation \cite{mckinsey2025stateai}. In parallel, Deloitte reports that GenAI is already being used in retail and consumer products for marketing and branding, audience segmentation, copy generation, sentiment analysis, and hyper-personalised shopping support, even though most organisations in these sectors remain in the early stages of adoption \cite{deloitte2025retailgenai}. These figures are useful not as the main frame of the argument, but as evidence that AI-mediated communication has become a practical and visible feature of contemporary digital media environments.

The rapid diffusion of generative artificial intelligence (GenAI) has intensified this transformation. In contrast to earlier forms of predictive AI, which primarily supported classification, targeting, and forecasting, GenAI can directly produce marketing content such as text, images, product descriptions, promotional copy, and interactive brand messages. As a result, AI is no longer only involved in analysing consumer data; it is also participating in the creation of the marketing message itself \cite{hermann2024consumer,ubal2024framework}. This development raises new questions about whether consumers interpret AI-mediated communication in the same way as human-created communication.

From an interactive digital media perspective, this shift matters because consumer responses are not determined solely by functional efficiency. Consumers also evaluate whether an AI-mediated message, interface, or interaction seems trustworthy, authentic, relevant, credible, intrusive, or helpful. These judgments shape broader outcomes such as attitudes towards the brand, willingness to engage, and purchase intention. Although AI can improve convenience and personalisation, its use may simultaneously generate unease when consumers perceive automation as impersonal, deceptive, or overly dependent on personal data \cite{kumar2024aipowered,mohdamin2025adoption}. Understanding consumer perception is therefore central to understanding not only AI adoption in marketing, but also the changing character of digital communication and mediated interaction.

\section{Research Problem}
The use of AI in marketing presents a tension between commercial opportunity and consumer acceptance. On the one hand, AI promises improved efficiency, scalability, and relevance. Personalised recommendations, automated customer support, and AI-generated promotional content can make marketing more responsive to consumer needs and organisational goals. On the other hand, these same practices may trigger concerns about privacy, transparency, manipulation, authenticity, and the erosion of human connection in brand communication \cite{ubal2024framework,hermann2024consumer}.

This tension becomes particularly visible in studies of AI-generated and AI-disclosed marketing content. Emerging empirical work shows that consumers often respond differently when they know a message, image, or advertisement has been created by AI rather than by a human. For example, emotional marketing communication attributed to AI can reduce perceived authenticity and weaken favourable consumer responses \cite{kirk2025aiauthorship}. Similarly, disclosure labels can increase novelty while also decreasing authenticity, producing mixed effects on attitudes and purchase intention \cite{shi2026disclosure}. Studies on AI-generated advertising and visual content also suggest that trust and purchase intention may decline when AI authorship is salient, especially in high-involvement or emotionally sensitive contexts \cite{israfilzade2025advertising,belanche2025images,zhang2025images}.

At the same time, not all findings are negative. Research on AI-powered personalisation indicates that relevance, usefulness, and trust can mediate positive consumer responses, particularly when AI applications are perceived as valuable rather than merely automated \cite{an2025personalized,srivastava2026social}. This suggests that consumer perception of AI marketing is neither uniformly positive nor uniformly negative. Instead, it appears to depend on a set of interconnected perceptual dimensions that operate differently across contexts, content types, and degrees of AI visibility.

Despite growing interest in this area, the literature remains fragmented. Existing studies often examine trust, authenticity, privacy, usefulness, or purchase intention as separate issues, and many focus on specific applications such as AI-generated advertising, recommendation systems, or disclosure labels. As a result, there is still a need for a more integrated account of how AI and GenAI influence consumer perception across marketing contexts. This dissertation addresses that need by synthesising the literature through a thematic review focused explicitly on the consumer perspective.

\section{Research Aim}
The aim of this dissertation is to examine how artificial intelligence and generative artificial intelligence shape consumer perception in AI-mediated communication across digital marketing contexts.

\section{Research Objectives}
\begin{enumerate}
    \item To review academic literature on AI and generative AI in digital marketing and brand communication.
    \item To identify the main consumer perception themes that recur across this literature.
    \item To analyse the mechanisms through which trust, authenticity, personalisation, privacy concerns, and perceived usefulness shape consumer responses to AI-mediated communication.
    \item To explain the conditions under which AI-mediated communication is more likely to generate favourable or unfavourable consumer responses.
\end{enumerate}

\section{Research Questions}
\begin{enumerate}
    \item What are the main dimensions through which AI and generative AI shape consumer perception in digital marketing communication?
    \item Through what mechanisms do trust, authenticity, personalisation, privacy concerns, and perceived usefulness influence consumer responses to AI-mediated communication?
    \item Under what conditions are consumer responses to AI and generative AI in digital marketing communication more likely to become positive or negative?
\end{enumerate}

\section{Significance of the Study}
This study is significant for both academic and practical reasons. Academically, it contributes to a developing body of literature at the intersection of AI, consumer behaviour, digital communication, and interactive media. While research on AI in marketing has expanded rapidly, the consumer perspective remains conceptually dispersed across different applications and theoretical lenses \cite{jain2024twodecades,ubal2024framework}. By bringing together literature on predictive AI, generative AI, advertising, personalisation, disclosure, and mediated interaction, this dissertation offers a clearer account of the perceptual mechanisms that shape consumer responses.

Practically, the study is relevant for marketers and organisations that are increasingly incorporating AI into customer-facing activities. Businesses may assume that efficiency and innovation automatically improve consumer evaluations, yet the literature suggests that acceptance depends on how consumers interpret the meaning of AI use. If AI-generated content is seen as useful and relevant, it may strengthen engagement and purchase intention; if it is seen as inauthentic, misleading, or invasive, it may damage trust and brand relationships \cite{kirk2025aiauthorship,an2025personalized,israfilzade2025advertising}. A better understanding of these dynamics can help organisations deploy AI more responsibly and strategically.

The study is also relevant to interactive digital media design and digital communication practice. As AI becomes embedded in chat interfaces, recommendation systems, disclosure cues, automated service encounters, and generative content workflows, design decisions increasingly shape how AI is experienced by users. The dissertation therefore has implications for how AI-mediated brand interactions are framed, how disclosure is designed at the user experience level, and how organisations position GenAI content within broader digital media production and communication processes.

\section{Scope of the Study}
This dissertation focuses on consumer perception in digital marketing communication rather than on the technical performance of AI systems. It examines literature on AI-enabled practices including personalised advertising, recommendation systems, AI-generated advertising content, chatbots, and social media communication. The study considers both traditional predictive AI and newer forms of generative AI, especially where AI becomes visible in the marketing message or consumer experience.

More specifically, the dissertation is concerned with consumer-AI interaction in digital media environments rather than with organisational adoption at the level of marketing strategy alone.

The dissertation is limited to secondary research and does not involve primary data collection. Its purpose is not to test causal relationships empirically, but to synthesise and interpret existing academic evidence. The emphasis is therefore on identifying recurring themes, tensions, and patterns in the literature rather than on evaluating the engineering performance of AI technologies.

\section{Dissertation Structure}
This dissertation is organised into five chapters. Chapter 1 introduces the research background, problem, aim, objectives, research questions, and scope of the study. Chapter 2 reviews the literature on AI, generative AI, and consumer perception in digital marketing communication, and develops the conceptual basis for the study. Chapter 3 explains the methodology, including the literature selection process and thematic analysis approach. Chapter 4 presents and discusses the main themes identified from the reviewed literature. Chapter 5 concludes the dissertation by summarising the findings, outlining the implications of the study, and suggesting directions for future research.

\chapter{Literature Review}

\section{Introduction}

This chapter reviews existing academic literature on artificial intelligence (AI) and generative AI in marketing, with a particular focus on consumer perception. It aims to identify key themes and theoretical perspectives that explain how AI influences consumer attitudes and responses.

\section{Artificial Intelligence in Marketing}

Artificial intelligence has become an important tool in marketing. It is widely used in recommendation systems, personalised advertising, chatbots, and customer analytics. AI allows firms to process large amounts of data and deliver more targeted and efficient marketing strategies.

However, the use of AI also changes how consumers interact with marketing content. Instead of human-driven communication, many interactions are now mediated by algorithms and automated systems.

\section{Generative AI in Marketing}

Generative AI refers to technologies that can create content such as text, images, and videos. In marketing, it is used to generate advertisements, product descriptions, and social media content.

Compared to traditional AI, generative AI plays a more direct role in shaping marketing communication. This raises new questions about authenticity, creativity, and consumer trust.

\section{Consumer Perception}

Consumer perception refers to how individuals interpret and evaluate marketing information. It includes cognitive and emotional responses such as trust, perceived quality, credibility, and attractiveness.

In AI-driven marketing, perception becomes more complex because consumers may not always know whether content is created by humans or machines.

\section{Key Themes in AI Marketing and Consumer Perception}

Based on existing literature, several key themes can be identified:

\subsection{Trust}

Trust is one of the most important factors in AI marketing. Consumers may trust AI systems for their efficiency, but may also question their transparency and intentions.

\subsection{Authenticity}

AI-generated content can reduce perceived authenticity, especially when emotional or human-like communication is involved.

\subsection{Personalisation}

AI enables highly personalised marketing, which can improve relevance and satisfaction. However, excessive personalisation may also raise concerns.

\subsection{Privacy Concerns}

Consumers may feel uncomfortable when their data is used for AI-driven marketing. Privacy concerns can negatively affect attitudes and trust.

\subsection{Purchase Intention}

AI influences consumer decision-making, especially through recommendations and targeted advertising.

\section{Research Gap}

Although existing studies have explored AI in marketing, there are still gaps in understanding how generative AI specifically affects consumer perception. In particular, more research is needed to examine the interaction between trust, authenticity, and personalisation.

\section{Summary}

This chapter has reviewed key literature on AI and generative AI in marketing. It has identified major themes related to consumer perception, which will be further analysed in the methodology and findings chapters.
\chapter{Methodology}

\section{Research Approach}

This study adopts a qualitative research approach based on a literature review. Rather than collecting primary data, the research focuses on analysing existing academic studies related to artificial intelligence (AI), generative AI, and consumer perception in marketing.

\section{Research Design}

The study uses thematic analysis to identify patterns and themes across selected academic literature. This approach allows the research to organise and interpret complex findings from multiple sources.

\section{Data Collection}

The data for this study consists of academic journal articles related to AI in marketing and consumer perception. These articles were selected using academic databases such as Google Scholar, Scopus, and ScienceDirect.

Keywords used for searching include:

\begin{itemize}
    \item Artificial intelligence in marketing
    \item Generative AI advertising
    \item Consumer perception AI
    \item AI-generated content
    \item AI trust, authenticity, and privacy
\end{itemize}

\section{Data Selection Criteria}

To ensure relevance and quality, the following criteria were applied:

\begin{itemize}
    \item Peer-reviewed journal articles
    \item Published between 2020 and 2025
    \item Focus on AI or generative AI in marketing
    \item Discuss consumer perception or related concepts
\end{itemize}

A total of approximately 15–25 articles were selected for detailed analysis.

\section{Thematic Analysis}

Thematic analysis was used to identify key themes across the selected literature. The process followed several steps:

\begin{enumerate}
    \item Familiarisation with the data through reading the articles
    \item Initial coding of key concepts and findings
    \item Grouping codes into broader themes
    \item Reviewing and refining themes
\end{enumerate}

Key themes identified include trust, authenticity, personalisation, privacy concerns, and purchase intention.

\section{Limitations}

This study has several limitations. First, it relies only on secondary data, which may limit the depth of analysis. Second, the selection of literature may introduce bias. Third, the findings are based on interpretation and may not represent all consumer perspectives.

\section{Summary}

This chapter has outlined the research approach, data collection process, and analytical method used in this study. The next chapter presents the findings based on thematic analysis.
\chapter{Findings and Discussion}

\section{Introduction}
This chapter presents the findings from the thematic analysis of the selected literature and discusses their implications for understanding consumer perception in AI-driven marketing. Rather than treating positive and negative outcomes as contradictory findings, the discussion interprets them as recurring trade-offs and conditional effects. Following Puntoni et al.'s experiential framework \cite{puntoni2021experiential}, the chapter is organised around four recurring tensions: data capture, classification, delegation, and social interaction. This structure makes it possible to show not only what patterns appear in the literature, but also why those patterns arise and under what conditions consumer responses become more favourable or more resistant.

For clarity, the analysis distinguishes between three levels that are closely related but not identical: interpretive appraisal, affective response, and behavioural outcome. In the chapter, consumer perception refers most directly to the first two levels, that is, how consumers make sense of AI and how they feel about it. Behavioural outcomes such as engagement, purchase intention, avoidance, or resistance are treated as consequences shaped by those perceptions rather than as perception itself.

\section{Interpreting the Four Tensions}
The four thematic groups identified in this dissertation should not be understood as four separate lists of positive and negative findings. Rather, each of them captures a recurring tension in how consumers interpret AI in marketing. These tensions help explain why consumer perception of AI is rarely stable or one-dimensional. Across the literature, consumers often value what AI provides while simultaneously resisting the conditions through which those benefits are delivered.

The first tension is between being served and being exploited. AI-driven marketing can make the consumer experience more relevant, efficient, and convenient, but these benefits depend on extensive data capture. As a result, consumers are often drawn into a privacy calculus in which functional value is weighed against surveillance, loss of control, and mistrust. The second tension is between being understood and being misunderstood. AI classification and personalisation can make consumers feel recognised, yet the same systems may also make them feel stereotyped, reduced to a category, or stripped of individuality.

The third tension is between being empowered and being replaced. Delegating tasks to AI can reduce effort and improve efficiency, but it can also undermine autonomy, self-efficacy, and the sense of authorship over one's own choices. The fourth tension is between being connected and being alienated. AI-mediated communication can create speed, responsiveness, and even novelty, but it may also feel emotionally empty or morally inappropriate when consumers expect genuine empathy, warmth, or human sincerity.

Taken together, these tensions suggest that the literature does not simply contain contradictory findings. Instead, it shows that consumer responses to AI are shaped by recurring trade-offs. AI is more likely to be perceived positively when it delivers practical value without violating expectations of privacy, individuality, autonomy, or authenticity. It is more likely to be perceived negatively when those benefits come at too high a psychological, emotional, or ethical cost. The four sections that follow develop this argument in greater detail.

\section{Overview of the Coding Framework}
The review suggests that consumer perception of AI is shaped by four recurring tensions:

\begin{itemize}
    \item Data capture: served versus exploited
    \item Classification: understood versus misunderstood
    \item Delegation: empowered versus replaced
    \item Social interaction: connected versus alienated
\end{itemize}

These tensions do not operate independently. Across the literature, consumer responses to AI are shaped by the interaction between perceived benefits such as convenience, relevance, efficiency, and responsiveness, and perceived costs such as intrusiveness, opacity, stereotyping, loss of autonomy, and perceived inauthenticity. The findings therefore suggest that AI in marketing should be understood through a relational and interpretive lens, not simply through a performance lens. Trust, authenticity, privacy concern, and behavioural intention are best understood as outcomes or mechanisms distributed across these four broader experience domains rather than as fully separate themes.

\section{Data Capture and Privacy: Served versus Exploited}
The first thematic tension concerns data capture. AI-driven marketing depends heavily on consumer data in order to personalise recommendations, optimise advertising, and support real-time interaction. From the consumer perspective, this can create a positive experience when data use produces relevance, convenience, and reduced search effort. In such cases, consumers may feel served by AI because the system appears to anticipate their needs and improve the efficiency of the market experience \cite{an2025personalized,ameen2021customer}.

However, the same process can produce a very different perception when data collection feels opaque, excessive, or difficult to control. Aguirre et al. show that personalised advertising performs better when information collection is overt, but that covert collection heightens perceived vulnerability and reduces effectiveness \cite{aguirre2015paradox}. Hardcastle et al. similarly show that AI-driven personalised journeys can create frustration and resentment when consumers feel persistently monitored or when stored data seems overly extensive \cite{hardcastle2025journeys}. Boerman et al. likewise show that targeted digital advertising can quickly shift from relevance to perceived intrusiveness when consumers become more aware of how data is being collected and used \cite{boerman2017oba}. Martin and Murphy provide broader conceptual support by showing that privacy in marketing is not a peripheral issue but a structural concern tied to trust, legitimacy, and power in data use \cite{martin2017privacy}. More broadly, Davenport et al. and Kumar et al. identify privacy and data governance as structural tensions in AI-enabled marketing \cite{davenport2020future,kumar2024aipowered}. Ozturkcan and Bozda{\u{g}} extend this concern by showing how opaque practices and exaggerated claims about AI can generate a broader cycle of mistrust, in which AI washing is followed by consumer backlash or ``AI booing'' \cite{ozturkcan2025mistrust}. Srivastava and Gurme further show that privacy concern and perceived surveillance can directly weaken trust even when personalisation remains technically effective \cite{srivastava2026social}. These studies suggest that the issue is not simply whether firms use data, but whether that use is interpreted as proportionate, transparent, and legitimate.

What consumers perceive in this theme is a tension between being efficiently served and being quietly exploited. Consumers value tailored relevance, convenience, and reduced search effort, but they also become highly sensitive to signs that their personal data is being continuously monitored, stored, or abstracted into computational profiles without meaningful consent. This is why the same AI system can be interpreted either as useful service or as invasive surveillance.

Why this perception occurs can be explained through privacy calculus and perceived vulnerability. Consumers weigh the practical benefits of personalisation against the costs of dataveillance, loss of privacy, and reduced control. This tension becomes especially intense when data collection is covert or when the logic of the system remains opaque. It is further exacerbated when firms overstate their AI sophistication or ethical responsibility, because AI washing inflates expectations that are then violated when the system appears exploitative rather than trustworthy. The underlying mechanism is therefore not merely data use itself, but the interpretation of that data use as either legitimate service or covert extraction.

The main behavioural outcomes affected by this tension are trust, click-through intention, ad avoidance, purchase intention, and longer-term willingness to remain engaged with the customer journey. Transparent and proportionate data practices are more likely to sustain trust and purchase intention, whereas opaque or excessive data capture is more likely to produce avoidance, reactance, and, in more severe cases, public backlash or disengagement \cite{an2025personalized,srivastava2026social}. The evidence for this theme is especially strong because similar patterns appear across experiments, surveys, and review-based studies. The tension is therefore best understood not as a contradiction in the literature, but as a privacy-personalisation trade-off built into the data capture experience itself.

\section{Classification and Personalisation: Understood versus Misunderstood}
The second thematic tension concerns classification. AI systems sort consumers into categories, infer their likely preferences, and use those classifications to generate personalisation. When this process works well, consumers may feel understood because the message, recommendation, or service appears relevant to their actual needs. This is why personalisation is often associated with favourable consumer response in the literature \cite{an2025personalized,haleem2022ai,kumar2024aipowered}.

Yet the classification process can also make consumers feel misunderstood. Puntoni et al. suggest that AI can produce experiences of misunderstanding when it reduces people to narrow inferred categories or makes judgements that fail to reflect their full identity and intention \cite{puntoni2021experiential}. Hou et al. provide particularly direct evidence for this problem by showing that consumers often infer that AI-generated personalised solutions are shared with a larger group of people and are therefore less unique \cite{hou2026uniqueness}. Starke et al. further show that perceptions of algorithmic fairness are highly sensitive to context, procedure, and outcome, which helps explain why classification can also trigger concerns about stereotyping, bias, and unequal treatment \cite{starke2022fairness}. Longoni et al. add an important boundary condition by showing that consumers become especially resistant when AI is perceived to lack the judgement needed for high-stakes or personally nuanced decisions \cite{longoni2019medical}. Even where personalisation is technically accurate, it may still feel mechanical or impersonal if the consumer believes that many others are being treated in essentially the same way. The issue is therefore not only predictive success, but whether the classification feels sufficiently individual and context-sensitive.

What consumers perceive in this theme is a tension between being recognised and being reduced. On the positive side, AI classification can make consumers feel understood when the system appears to anticipate needs and provide relevant solutions. On the negative side, consumers may feel misunderstood when they suspect that the system relies on stereotypes, rigid templates, or overgeneralised profiles rather than on an appreciation of their individual complexity.

Why this perception occurs can be explained through uniqueness neglect and fairness concerns. Because consumers often attribute low experiential sensitivity to AI, they assume that algorithmic recommendations are generated through broad rules rather than genuine insight. This leads them to infer a larger expected group size, meaning that the personalised solution does not feel rare, exclusive, or truly individual. Where biased or historically skewed data are involved, this perception becomes more severe because classification can also be interpreted as unfair or discriminatory. The central mechanism is therefore not simply whether the recommendation is personalised, but whether the underlying classification appears psychologically sensitive and procedurally fair.

The behavioural outcomes affected by this theme include adoption intention, willingness to rely on AI advice, willingness to pay for personalised solutions, and acceptance of AI providers in higher-stakes contexts. When classification produces a convincing sense of understanding, AI is more likely to be accepted as useful and competent. When it produces uniqueness neglect or fairness concerns, consumers become more reluctant to trust the recommendation, less willing to adopt it, and more likely to reject AI providers in contexts where individual nuance matters. The evidence base here is moderate to strong: the overall pattern is clear, even though fairness perceptions remain highly context-sensitive. The classification experience therefore introduces a subtle but important trade-off between precision and reductionism.

\section{Task Delegation and Autonomy: Empowered versus Replaced}
The third thematic tension concerns delegation. Consumers often rely on AI to assist with search, recommendation, customer service, and decision support. In many contexts this delegation is experienced positively because it reduces effort, saves time, and makes the market environment easier to navigate. Ameen et al. show that AI-enabled services can improve customer experience when they are perceived as competent, convenient, and supportive of the consumer's goals \cite{ameen2021customer}. In these situations, delegation is experienced as empowerment rather than dependence.

At the same time, delegation can generate a loss of autonomy. Puntoni et al. argue that AI may create tension when consumers feel that judgement, choice, or effort has been transferred too far away from the self \cite{puntoni2021experiential}. Hardcastle et al. suggest that this tension intensifies when predictive journeys repeatedly direct consumers toward similar content or options, effectively narrowing what they see and reinforcing a sense of constrained choice \cite{hardcastle2025journeys}. Davenport et al. similarly note that AI-driven micro-targeting can reduce search costs while also undermining the consumer's perceived freedom of choice \cite{davenport2020future}. This becomes especially relevant when AI does not merely support a decision but appears to narrow the range of options, shape preferences too strongly, or replace forms of action that consumers associate with agency, creativity, or competence. In such cases, the consumer no longer experiences AI as assistance but as partial replacement.

What consumers perceive in this theme is a tension between empowerment and displacement. AI can feel empowering when it removes mundane effort and helps consumers navigate large amounts of information more efficiently. Yet the same delegation can feel threatening when consumers sense that control is shifting away from them, especially in tasks that involve subjectivity, judgement, or identity.

Why this perception occurs can be explained through perceived sacrifice and loss of autonomy. AI systems can guide choices so heavily that consumers no longer experience themselves as the true authors of their decisions. This becomes more serious when algorithmic curation narrows visible options or creates filter bubbles, because the convenience of delegation starts to look like a loss of freedom. In identity-relevant activities, delegation may also feel like the loss of an accomplishment that consumers believe should belong to them rather than to a machine. The core mechanism here is the conversion of convenience into psychological threat when support turns into substitution.

The behavioural outcomes associated with this theme include resistance to recommendations, frustration, reduced reliance on AI support, and psychological reactance. In some cases, consumers may deliberately ignore or reject algorithmic guidance simply to reassert their independence. Although the empirical base on long-term autonomy loss is still developing, the existing evidence consistently suggests that delegation is most positively perceived when it remains supportive and reversible, and most negatively perceived when it appears to replace rather than augment the consumer's own agency.

\section{Social Interaction and Authenticity: Connected versus Alienated}
The fourth thematic tension concerns social interaction. Consumers increasingly encounter AI in socially framed contexts, such as chatbots, AI-generated advertising, virtual assistants, and other forms of automated brand communication. These interactions can create a sense of connection when they are responsive, clear, and useful. Consumers may appreciate instant assistance, reduced waiting time, and the novelty of AI-mediated brand interaction, especially in factual or utilitarian settings \cite{ameen2021customer,chan2025genai}.

However, this same social experience can become alienating when AI communication feels emotionally empty, morally inappropriate, or insufficiently human. This is where authenticity becomes especially important. As discussed in the literature review, authenticity in this dissertation has two dimensions: source authenticity and emotional authenticity. Source authenticity concerns whether the message is perceived as genuinely human-authored or AI-authored \cite{kirk2025aiauthorship}. Emotional authenticity concerns whether the communication feels sincere, empathetic, and genuinely intended \cite{arango2023charitable}. These two dimensions help explain why consumers may resist AI not simply because it is artificial, but because it disrupts expected forms of warmth, empathy, or symbolic meaning.

What consumers perceive in this theme is a dual response: AI can feel novel, responsive, and efficient, yet also inauthentic, emotionally empty, or even deceptive. Consumers may enjoy the speed and availability of AI-mediated communication, but they remain aware that machines do not possess genuine internal emotional states. Kirk and Givi show that this can produce an AI-authorship effect in which the perceived lack of real feeling weakens response to emotionally expressive communication \cite{kirk2025aiauthorship}. Shi and Jiang similarly show that AI disclosure can increase novelty while simultaneously reducing authenticity \cite{shi2026disclosure}. This creates a tension between the convenience of connection and the discomfort of artificiality.

Why this perception occurs can be explained through two competing mechanisms. The first is perceived novelty, which can stimulate curiosity and engagement because AI signals innovation. The second, often more powerful, is perceived inauthenticity. When AI attempts to simulate empathy, warmth, or emotional sincerity, consumers may judge the interaction as fake because the expression is not backed by genuine feeling. Arango et al. show how this problem can reduce empathy and donation intention in charitable advertising \cite{arango2023charitable}, while Belanche et al. show that resistance becomes stronger in hedonic or high-involvement settings where realism and genuine human presence matter more \cite{belanche2025images}. In some contexts this mismatch generates moral disgust, especially when AI is used in emotionally charged or morally sensitive communication. The effect is therefore strongly context-dependent: utilitarian tasks invite more acceptance, whereas hedonic or high-involvement experiences intensify the need for human touch and authentic presence.

The behavioural outcomes associated with this theme include short-term engagement through novelty, but also reduced word-of-mouth, lower loyalty, weaker donation intention, and stronger resistance when authenticity concerns dominate. In medical, charitable, and emotionally expressive settings, alienation can become especially strong because consumers fear that their individuality, seriousness, or emotional needs are not being genuinely recognised \cite{longoni2019medical}. The evidence for this theme is strong across multiple experiments and applied contexts. Social interaction with AI is therefore best understood as a double-edged mechanism that can create both connection and alienation depending on how the communication is framed and where it occurs.

\section{How Trust, Authenticity, and Behavioural Intention Fit Across the Four Themes}
The analysis suggests that trust, authenticity, and behavioural intention should not be treated as isolated findings detached from the broader experience structure. In the literature review, trust was discussed as a distinct theme because it is one of the most visible concepts in prior studies and often appears as an explicit dependent or mediating variable. In the findings, however, trust is repositioned as a cross-cutting mechanism because the thematic synthesis shows that it is produced through several different experience domains rather than operating as one separate empirical topic. Trust is shaped across data capture, classification, delegation, and social interaction, which is why it is analytically more useful here as an evaluative mechanism than as a fifth standalone theme.

Authenticity is primarily concentrated in the social interaction theme, but it also interacts with trust because doubts about authorship and emotional sincerity can weaken the credibility of AI-mediated communication. Behavioural intention, by contrast, is best understood as an outcome of how the four tensions are resolved. Consumers are more likely to engage, purchase, or respond favourably when they feel served, understood, empowered, and connected. They are less likely to do so when they feel exploited, misunderstood, replaced, or alienated.

\section{How the Four Tensions Interact}
The four tensions identified in this dissertation should not be treated as independent blocks. Although they are presented separately for analytical clarity, the literature suggests that they often reinforce or condition one another. In practice, consumers do not encounter data capture, classification, delegation, and social interaction as isolated dimensions. They encounter them together within the same digital experience, which means that one tension can intensify or moderate another.

The clearest interaction appears between data capture and classification. Consumers are less likely to feel understood by personalisation when they are already uneasy about how their data was collected. In this sense, privacy concern can weaken the positive effects of relevance, because the same recommendation may be interpreted less as helpful recognition and more as intrusive profiling. A similar interaction exists between classification and delegation. When AI appears to rely on narrow categories or generic rules, consumers are less willing to delegate judgement to it, especially in contexts that require nuance, taste, or personal identity. Misunderstanding therefore makes replacement feel more threatening.

Data capture also interacts with social interaction. AI-mediated communication may already feel emotionally thin or inauthentic, but that concern can become stronger when the message is also understood as the product of hidden surveillance or excessive profiling. In other words, perceived intrusion can make perceived inauthenticity more severe. Conversely, when AI use is transparent, proportionate, and clearly bounded, consumers may be more willing to accept novelty or limited automation in communication.

Taken together, these patterns suggest that the four tensions form an interpretive system rather than a simple list. Data concerns can shape whether personalisation feels caring or exploitative. Classification can shape whether delegation feels supportive or substitutive. Social interaction can amplify how questions of trust, privacy, and authenticity are ultimately resolved. This interactional reading strengthens the framework by showing that consumer perception is shaped not only by the presence of individual tensions, but by how those tensions accumulate and condition each other within a single AI-mediated encounter.

\section{Conditions That Shape More Positive or Negative Responses}
The literature suggests that consumer responses to AI become more positive or more negative under identifiable conditions. One important condition is the visibility of AI. When AI works in the background to improve convenience or targeting, it is often evaluated more positively. When it becomes highly visible as the creator of emotionally meaningful or identity-relevant communication, consumers become more alert to questions of authenticity and trust.

Another condition is message type. AI appears to be more favourably received when the message is informational, functional, or service-oriented than when it is expected to convey emotional sincerity, artistic originality, or social meaning. Product and context also matter. High-involvement, hedonic, and symbolic categories tend to intensify concerns about authenticity and human judgement, whereas lower-stakes or convenience-driven contexts often make efficiency benefits more salient \cite{belanche2025images,arango2023charitable}.

Data practice is a further condition. Consumers are more willing to accept AI when personalisation feels proportionate and understandable. They become more resistant when the same personalisation feels opaque, excessive, or difficult to control \cite{boerman2017oba,srivastava2026social}. These patterns suggest that the consumer response to AI is not random. It is shaped by recurring contextual conditions that determine whether the benefits of AI outweigh its relational and ethical costs.

\section{Implications of the Findings}
For researchers, the findings suggest that future work should move beyond simply identifying whether AI has positive or negative effects. What matters more is specifying the mechanisms and conditions through which those effects occur. This would allow future studies to build a more cumulative understanding of consumer response to AI and GenAI across different marketing settings.

For practitioners, the main implication is that AI should not be managed only as an efficiency tool. It must also be managed as a communicative and relational presence. The success of AI in marketing depends on whether it is used in ways that preserve trust, respect privacy expectations, and match the symbolic demands of the message context. In practice, this means that high precision, fast content generation, or automated interaction are not enough on their own. They must be aligned with what consumers interpret as appropriate, credible, and worth engaging with.

\section{Summary}
This chapter has shown that consumer perception of AI in marketing is shaped by recurring trade-offs rather than by simple contradiction. Organised through the four tensions of data capture, classification, delegation, and social interaction, the analysis suggests that AI tends to generate more positive responses when consumers feel served, understood, empowered, and connected, but more negative responses when they feel exploited, misunderstood, replaced, or alienated. These outcomes are produced through identifiable mechanisms and shaped by identifiable conditions, including AI visibility, message type, emotional expectation, and data practice. The next chapter concludes the dissertation by summarising these findings and outlining their broader implications.

\chapter{Conclusion}

\section{Introduction}

This chapter summarises the key findings of the study, discusses theoretical and practical implications, and outlines limitations and directions for future research.

\section{Summary of Findings}

This dissertation examined how artificial intelligence (AI) and generative AI influence consumer perception in marketing. Through a thematic analysis of existing literature, five key themes were identified: trust, authenticity, personalisation, privacy concerns, and purchase intention.

The findings suggest that AI has a dual impact on consumer perception. On one hand, AI improves efficiency, relevance, and personalisation in marketing. On the other hand, it raises concerns related to authenticity, transparency, and privacy. These conflicting effects indicate that AI-driven marketing involves a trade-off between benefits and risks.

Trust was found to be a central factor linking all other themes. Consumers are more likely to accept AI-driven marketing when they perceive it as trustworthy and useful. However, when authenticity is questioned or privacy concerns arise, trust may decrease, leading to lower purchase intention.

\section{Theoretical Implications}

This study contributes to the literature by integrating research on AI in marketing with consumer perception theories. It highlights the importance of examining multiple interconnected factors rather than focusing on a single variable.

In particular, the study emphasises the role of authenticity and trust in shaping consumer responses to AI-generated content. It also extends existing research by considering generative AI as a distinct category that directly influences marketing communication.

\section{Practical Implications}

From a practical perspective, the findings suggest that marketers should carefully balance the use of AI technologies with consumer expectations.

First, transparency is important. Clearly communicating the use of AI may help reduce uncertainty and increase trust.

Second, firms should ensure that AI-generated content maintains a level of authenticity, especially in emotional or human-centred communication.

Third, data usage should be managed responsibly to reduce privacy concerns and build long-term consumer trust.

\section{Limitations}

This study has several limitations. It relies on secondary data and does not include primary empirical research. The selection of literature may also introduce bias, as only a limited number of studies were analysed.

In addition, the findings are based on interpretation and may not fully capture the diversity of consumer experiences across different cultural and market contexts.

\section{Future Research}

Future research could address these limitations by conducting empirical studies, such as surveys or experiments, to test the relationships identified in this study.

Further research is also needed to explore how different types of AI applications influence consumer perception in various contexts. In particular, the impact of generative AI on authenticity and trust remains an important area for investigation.

\section{Final Remarks}

Artificial intelligence and generative AI are transforming the marketing landscape. While they offer significant opportunities for improving efficiency and personalisation, they also introduce new challenges for consumer perception.

Understanding how consumers respond to AI-driven marketing is therefore essential for both researchers and practitioners. This study provides a foundation for future research in this evolving field.

% note that your supervisor may have a strong opinion on the style of referencing you use. Some background is available at https://www.overleaf.com/learn/latex/Bibtex_bibliography_styles
\bibliographystyle{IEEEtran} %Changed to IEEETran by HS
%\bibliographystyle{unsrt}

\bibliography{bibs/ai_marketing_intro}
\appendix
\renewcommand{\thechapter}{A\arabic{chapter}}
\chapter{Appendix}

\section{Coding Framework Overview}
This appendix provides a summary of the thematic coding framework used in the dissertation. The coding process was informed by Puntoni et al.'s experiential framework and translated into four overarching consumer-perception tensions. The appendix is included to show how the main analytical categories used in Chapter 3 and Chapter 4 were derived from the literature.

\begin{table}[h!]
\centering
\caption{Summary of the thematic coding framework}
\begin{tabular}{p{3cm}p{3.2cm}p{3.3cm}p{4cm}}
\toprule
\textbf{Theme group} & \textbf{Puntoni domain} & \textbf{Core tension} & \textbf{Illustrative constructs} \\
\midrule
Data capture and privacy & Data capture experience & Served versus exploited & Privacy calculus, dataveillance, perceived vulnerability, transparency, trust \\
\midrule
Classification and personalisation & Classification experience & Understood versus misunderstood & Personalisation, uniqueness neglect, expected group size, stereotyping, algorithmic fairness \\
\midrule
Task delegation and autonomy & Delegation experience & Empowered versus replaced & Self-efficacy, perceived control, perceived sacrifice, autonomy, algorithm aversion \\
\midrule
Social interaction and authenticity & Social experience & Connected versus alienated & Source authenticity, emotional authenticity, empathy gap, moral disgust, novelty \\
\bottomrule
\end{tabular}
\end{table}

\section{Coding Table}
Table A1 provides a fuller version of the coding logic used in the dissertation. It links the main theme groups to associated concepts, supporting studies, and their interpretive relevance to consumer perception.

\begin{table}[h!]
\centering
\caption{Full thematic coding table}
\begin{tabular}{p{2.8cm}p{3.4cm}p{4.4cm}p{4cm}}
\toprule
\textbf{Theme group} & \textbf{Key constructs} & \textbf{Illustrative supporting studies} & \textbf{Interpretive relevance} \\
\midrule
Data capture and privacy & Privacy calculus, surveillance, dataveillance, trust, transparency, perceived exploitation & Aguirre et al. (2015); Boerman et al. (2017); Martin and Murphy (2017); Hardcastle et al. (2025); An and Ngo (2025); Srivastava and Gurme (2026); Ozturkcan and Bozda{\u{g}} (2025) & Shows how consumers may value tailored service while simultaneously resisting opaque or excessive data capture \\
\midrule
Classification and personalisation & Personalisation, algorithmic classification, uniqueness neglect, expected group size, bias, fairness & Puntoni et al. (2021); Hou et al. (2026); Starke et al. (2022) & Shows that AI may create relevance without necessarily creating a feeling of individuality or fair recognition \\
\midrule
Task delegation and autonomy & Delegation, self-efficacy, autonomy, perceived control, replacement, psychological threat & Puntoni et al. (2021); Davenport et al. (2020); Ameen et al. (2021); Hardcastle et al. (2025) & Shows how AI can reduce effort while also narrowing choice and weakening the consumer's sense of agency \\
\midrule
Social interaction and authenticity & Source authenticity, emotional authenticity, empathy gap, novelty, alienation, moral disgust & Kirk and Givi (2025); Shi and Jiang (2026); Arango et al. (2023); Belanche et al. (2025); Longoni et al. (2019) & Shows how AI-mediated communication may be accepted in factual contexts but resisted in emotional or high-involvement settings \\
\bottomrule
\end{tabular}
\end{table}

\section{Appendix Note}
The appendix coding tables are included to make the analytical process transparent. They are not intended to replace the discussion in the main chapters, but to document how the dissertation moved from literature selection to thematic interpretation.



\end{document}
